\subsection{开放式评估：人类、Arena 与 LLM Judge 的三角关系}

slides 54--56 转向 ChatGPT 一类开放式系统。答案空间不再有限，任务也可能同时涉及事实、帮助性、风格和安全。InstructGPT 使用成对人类偏好；Chatbot Arena 让真实用户在盲测中选择；AlpacaEval 用强 LLM 代替大量人工判断，以极低成本获得与 Arena 高相关的 win rate。这个进步解决了速度问题，却没有自动解决偏差问题，尤其是 judge 对长度、格式、自信语气或特定模型风格的偏好。

\lecturefigure{slides-images/slide-054.png}{官方 slide 54：开放式 aligned LLM 需要偏好评估}
\lecturefigure{slides-images/slide-055.png}{官方 slide 55：Chatbot Arena 的盲测与真实用户偏好}
\lecturefigure{slides-images/slide-056.png}{官方 slide 56：AlpacaEval 的 LLM judge、速度与 spurious correlation}

\begin{knowledgebox}{读图：相关性高不等于评估器无偏}
slide 56 的“98\% correlation”表达在特定模型集合和协议上，AlpacaEval 排序与 Arena 高度一致。相关性不能证明单条判断正确，也不能保证分布变化后仍成立。若候选模型学会输出更长答案，而 judge 恰好偏好长度，win rate 会提高但真实效用未必提高。解决方法包括 length-controlled regression、分桶比较、交换答案位置、多 judge 复核和定期用人类样本校准。
\end{knowledgebox}

\begin{importantbox}{开放式评估的最小证据包}
一个可用于研发决策的 open-ended 结果，至少应报告：任务样本来源、对比 baseline、judge 提示与版本、位置随机化、长度分布、胜负/平局规则、置信区间和人类校准子集。Agent 还需补充终局成功、过程违规、工具副作用与每次成功成本。只有“平均 judge 分”无法区分模型更会完成任务，还是更会取悦 judge。
\end{importantbox}

\textbf{课堂提示：}讲者将 evaluation 称为实践瓶颈，是因为数据和奖励最终都依赖它。一个系统若用同一 LLM 同时生成训练数据、筛选样本和充当最终 judge，误差可能形成闭环自我强化。工程上应尽量让 generator、verifier 与 audit source 多样化，并保留真实用户或可执行环境作为外部锚点。

